\section{Project Retrospective}
\label{chap:retrospective}
In retrospect, we recognize that many things in the project that went well, while others could have been done better. This chapter will reflect on what worked as expected and what did not. We will evaluate six essential aspects of the project: \textit{The prestudy}, \textit{technological choices}, \textit{development process}, \textit{testing phase}, \textit{results}, and \textit{our collaboration}.

\textbf{Prestudy}\\
The prestudy part of this master’s thesis builds on the findings from our specialization project. As for almost any research, a thorough prestudy can be remarkably valuable, as it lays the foundation for which new ideas and areas of study can emerge. Our specialization project gave us invaluable knowledge on what alternative games are, leading to the development of our \textit{alternative game framework}. The knowledge and framework were particularly helpful for the creative process, as it led to the game concept \textit{Hover}, without forcing the creative process. 

\textbf{Technological choices}\\
During the specialization project, we studied technologies that could be used to develop Hover. This study allowed us to make an informed choice for our technology stack. We used this knowledge as evaluation criteria, in addition to previous experience and the framework’s maturity.

Since we chose a cross-platform framework, we did not exclude any potential test participants. The data from our test population revealed that 63\% used Android and 37\% iOS. Thus, if we had chosen a native framework, it would probably be more challenging to recruit testers through our network. 

The technology stack we used ended up working as expected, with only minor hiccups. Most of the issues were related to the tracking (hovering) functionality in the game. It proved challenging to identify and test all edge cases that could occur while hovering due to varying factors such as the device, permission, software weaknesses, and simulator limitations.

\textbf{Development process}\\
The development phase includes several matters that we potentially could have executed better. However, we are satisfied with how it turned out. Though the team only consisted of two developers, we benefited greatly from using agile methodology because it enforced structure and order to the process. Using Trello was beneficial as it visualized the task, improved communication, and helped us identify and prioritize tasks. Completing tasks on the board also proved to be a motivational factor. 

Before we began developing, we agreed on code conventions and architecture. This decision was very beneficial as it made us focus on functionality without building up technical debt. We also used pull requests to ensure that we conveyed the conventions and the quality of code we were aiming for. Additionally, the code reviews helped us catch errors, find potential improvements, and make sure both of us understood the code.

Due to the Covid-19 situation, remote work was sometimes necessary. In these situations, consistent use of digital collaboration tools proved its value and ensured that the development continued smoothly and efficiently.

\textbf{Testing phase}\\
The ongoing pandemic also affected the testing phase. It did bring a lot of uncertainty leading up to and during testing. We considered pushing the testing back and also extending it to ensure that we had enough data for the results. However, the uncertainty could be confusing for the participants, so we were careful to communicate transparently. We also thought it would be beneficial to stick to the initial testing plan for the sake of the participants’ satisfaction and the project’s progress. Luckily, we felt that the two weeks of testing provided sufficient data to move forward with processing it. 

During testing, we discovered that some of the Android users experienced minor issues that abrupted their hovering experience. One of these issues caused the leaderboard to re-render every second, causing it to blink. Since the issues only appeared on some Android devices, it was difficult to reproduce and identify. From that, we learned that we should have tested the application more on different Android devices. 

\textbf{Results}\\
When we began processing the results, we discovered that we could have put more thought into the questionnaire. We found that some of the statements and questions could have been formulated better, removed, or added. In addition, one participant stated that they wished there was a \textit{not applicable} (N/A) option for the statements. Including this option might have resulted in different results.

We choose to focus more on the two groups in the results than what we originally planned. Consequently, the questionnaire could have been improved to acquire further insight into how the relationship between the participants affects the players’ engagement, motivation, enjoyment, and habits.

\textbf{Our Collaboration}\\
The collaboration between us has been solid and transparent throughout the project. Before we started, we discussed our motivation for the project and agreed on expectations for the collaboration. We also created a project plan that we followed and used Trello to verify our progress. The open communication helped us identify obstacles as they occurred. We also had weekly digital meetings with our supervisor, whose experience, guidance, and good ideas helped us through all project phases. 
